\section{Weekly Summary}

\begin{tcolorbox}[colback=yellow!10,colframe=black!50,title=AI-Generated Content]
\noindent The summary below was written by Claude (Anthropic) from that week's regression outputs and series descriptions. It has not been reviewed or fact-checked by a human, and should be read as an automated, informal recap -- not analysis.
\end{tcolorbox}

Welcome back to the weekly ritual where we take two economic series that have never met, force them into a linear regression, and then act astonished by whatever falls out. This week the random number generator served up wood product manufacturing in the Great Lakes, European cell phone subscriptions, discontinued wage-share statistics, Black Americans not in the labor force, intellectual property import charges, and the durable-goods slice of consumer spending. A truly harmonious guest list, all of whom got along about as well as you'd expect strangers to.

Monday's matchup pitted Great Lakes wood product GDP against mobile cellular subscriptions in Europe and Central Asia, which is exactly the kind of pairing this project exists to humiliate. The raw regression shrugged -- not even significant, a rare case of two series honestly admitting they have nothing to say to each other. But then, plot twist nobody asked for, the detrended version actually crossed the significance line. So after we strip out the trends, the leftover wiggles in Midwestern lumber output and the leftover wiggles in Eurasian phone plans quietly agree on something. Is it real? Almost certainly not. But it's the one result this week that got stronger after detrending, which technically makes it the most "interesting," in the way that finding a sock behind your dryer is technically an achievement.

Tuesday delivered the week's big drama queen: a discontinued measure of the wage-and-salary share of gross domestic income versus the number of Black Americans not in the labor force. The raw regression posted an R-squared of 0.89 and a p-value with more zeros than a lottery jackpot, the sort of number that makes you briefly believe you've discovered a Law of Economics. Then detrending walked in, calmly removed the shared drift these two long series had been coasting on, and the relationship deflated from "world-changing" to a much humbler still-significant-but-please-calm-down. It's the textbook trend mirage: most of that gorgeous fit was just two series wandering in the same general direction across decades, holding hands purely because time was passing.

Wednesday's pairing of intellectual property import charges against the durable-goods share of consumer spending did the same striptease -- a respectable raw R-squared of 0.32 and a significant p-value that promptly evaporated the instant we detrended, landing at a p-value of 0.44 that says, in no uncertain terms, "nothing to see here." As always, the standing reminder: these are randomly paired series, run through the bluntest possible statistics, with zero causal claims and even less dignity. Spurious correlation isn't the bug here, it's the entire personality of the project. See you next week for more of it.
